\section{관련 연구}\label{sec:related}

본 논문의 위치를 이해하려면 기존 활성 스티어링 방법들이 공유하는 \emph{특정한 빈자리}를 먼저 보는 것이 도움이 된다. SEKA~\citep{SEKA}와 AdaSEKA~\citep{AdaSEKA}는 contrastive question-answer 쌍에서 학습한 저랭크 사영을 attention key에 주입하여 단일 개념 편집 벤치마크에서 state-of-the-art 성능을 달성한다. ITI~\citep{ITI}와 CAA~\citep{CAA}는 residual stream에 activation direction을 더해 행동 분포를 바꾼다. PASTA~\citep{PASTA}는 특정 토큰의 post-softmax attention weight를 직접 재조정하고, Focus Directions~\citep{FocusDirections}는 key와 query를 함께 조절해 관련 토큰에 집중하도록 유도한다. 이 방법들의 공통된 특징은 모두 추가 finetuning 없이 추론 시에만 개입한다는 점이며, 그 점이 이들이 매력적인 도구인 이유다. 그러나 같은 공통성 위에 또 하나의 공통성이 얹혀 있다. \emph{어느 방법도 ``지금까지 무엇을 방출했는지''에 대한 명시적 상태를 담지 않는다.} 정지 key-side 증폭, residual direction, post-softmax 재조정, focus direction 모두 디코딩 단계와 무관한 고정된 변환이다. 정리~\ref{thm:no-memory}은 이 공통된 성질이 왜 멀티-툴 과제에서 한계가 되는지를 기계론적으로 고정한다. 본 논문은 이 기존 방법들을 비판하기 위해서가 아니라, 그들이 차지하지 않은 설계 공간 — 부호 반전된 query 축소 — 을 분명히 드러내기 위해 이 사실을 사용한다.

본 논문이 기존 흐름과 갈라지는 지점은 세 가지이며 하나의 근본적 선택에서 파생된다: ``key를 정지 증폭하지 말고 query를 \emph{부호가 있는} 회전으로 조작하라''. SEKA 계열은 $\beta>0$(amplification)만 탐색하고 Q-coverage 계열은 $\beta<0$(suppression)만 탐색했는데, 본 논문은 두 부호 모두 서로 다른 과제에서 이긴다는 것을 측정하고 그 부호를 사전 예측하는 정리를 제시한다. 첫째, 과제가 달라진다. 단일 개념 편집이나 prompt highlighting이 아니라, 한 질의에서 두 개 이상의 도구를 순차적으로 복원해야 하는 multi-tool selection을 직접 평가한다. 이 과제에서는 ``이미 선택한 방향을 다시 증폭한다''는 행동이 단지 중립적이 아니라 적극적으로 해롭다. 둘째, 연산자의 위치와 부호가 달라진다. key를 더하는 대신 query에서 빼고, 양의 계수 대신 음의 계수를 쓴다. 음의 Q-side projection은 이미 지배적인 온톨로지 성분을 눌러 거친 coverage prior로 작동하며, 정리~\ref{thm:q-sign}이 이 기하를 정확히 고정한다. 셋째, 기저의 출처가 달라진다. SEKA는 contrastive QA 쌍에서 $P_{\mathrm{pos}}$를 유도하지만, 본 논문은 도구 카탈로그의 4-facet 온톨로지를 Gram-Schmidt로 직교화해 $B_{\mathrm{ont}}$를 얻는다. 이 세 차이는 각각 독립적 기여가 아니라 ``부호를 뒤집어 query에서 빼라''는 한 선택의 세 측면이다.

벤치마크 선택은 이 서사를 지원하도록 좁게 이루어진다. MetaTool Subtask4~\citep{MetaTool}은 각 질의에 정확히 두 개의 정답 도구를 지정한 497개의 프롬프트로 구성되어, 순차 coverage를 가장 깨끗하게 측정할 수 있는 벤치마크다. $\tau^2$-bench retail~\citep{tau2bench}은 정답 도구 수가 0에서 13까지 분포하여 ``짧은 시퀀스''와 ``긴 시퀀스'' 체제가 같은 벤치마크 안에 공존한다. 이것이 regime-split을 측정 가능하게 만든다. BFCL~\citep{BFCL}의 parallel/multiple subset은 병렬 도구 호출의 공식 벤치마크로 교차 검증에 쓰인다. KV-cache compression~\citep{KIVI}은 본 논문의 초점이 아니며, 모든 실험은 Qwen2.5-Instruct 계열(주 실험은 7B, 보조 스윕은 1.5B/3B/14B)~\citep{qwen25}로 제한한다.
